\section{Open Questions}

The following five questions remain genuinely open after this replication.
Each is stated with its basis in the current work and a concrete next step.

\begin{enumerate}

\item \textbf{Mixed-state QAE.}
How does the QAE performance degrade when the input ensemble consists of mixed
states (density matrices) rather than pure states? The paper's formalism is
stated for pure states; extension to mixed inputs requires either purification
(doubling qubit count) or a redesigned cost, and the compression/fidelity
trade-off in that regime is not characterized.
\emph{Next step:} purify a controlled family of rank-$r$ mixed states on $n=4$
qubits into 8 qubits, run the QAE with trash-fidelity cost on the purified
ensemble, and compare compression achieved vs the classical rank $r$; also
implement the swap-test-based cost variant from the paper's Appendix and check
whether it handles mixed inputs more gracefully.

\item \textbf{Modern ansatz families.}
Can the hardware-efficient encoder ansatz be replaced with a modern variational
quantum ML architecture (data-reuploading, quantum-convolutional, or
equivariant circuits) to close the optimization gap seen at maximum compression
($k=1$ stuck at $F \approx 0.78$ in this replication)?
\emph{Next step:} reimplement with three alternative ansatze at fixed parameter
count, retrain each with both COBYLA and a gradient-based optimizer
(parameter-shift rule), and report the $k=1$ fidelity distribution over
restarts.

\item \textbf{Classical PCA / SVD baseline.}
How does the QAE compare quantitatively against classical PCA/SVD on the same
input state ensemble? For pure states in a low-rank subspace, classical SVD is
exact and free; a real advantage claim requires a regime where quantum
compression provably beats classical.
\emph{Next step:} for each replication test point, also compute the classical
rank-$k$ truncated-SVD reconstruction fidelity on the same state vectors in
$\mathbb{C}^{2^n}$; tabulate QAE-$F$ vs SVD-$F$ vs latent size; then design an
ensemble where quantum compression could plausibly beat the classical baseline
in resource terms.

\item \textbf{Noise robustness of the encoder.}
How robust is the trained encoder to realistic gate noise (depolarizing,
coherent over-rotation, readout) at levels typical of current superconducting
or trapped-ion hardware?
\emph{Next step:} rerun the trained encoders under a depolarizing channel with
per-gate error rates from $10^{-4}$ to $10^{-2}$ using a free simulator
(qiskit-aer or cirq), plot recon-fidelity vs noise strength for each latent
size, then attempt noise-aware retraining and quantify whether the encoder
becomes robust or training itself collapses.

\item \textbf{ML-driven latent-dimension selection.}
Can ML (meta-learning across ensembles, or a differentiable indicator of
effective rank) be used to select the latent qubit dimension $k$ automatically,
rather than requiring the user to pre-specify it or sweep by hand?
\emph{Next step:} implement a two-stage loop where an outer meta-optimizer
adjusts $k$ based on a soft penalty combining reconstruction fidelity and a
compression prior; test on ensembles of known effective rank (2-D, 4-D, full
rank) and check whether the automated $k$ selection matches the ground-truth
effective dimension within $\pm 1$ qubit.

\end{enumerate}
